\documentclass[10pt, a4paper]{article}
\usepackage{preamble}

\newcommand{\bnabla}[0]{\mbf{\nabla}}

\title{Mathematical Methods II \\
    \large Prereading Notes}
\author{Luke Phillips}
\date{October 2025}

\begin{document}

\maketitle

\newpage

\tableofcontents

\newpage

\section{Line integrals}

\subsection{Curves}
It is convenient to describe a curve by a parameterisation,
$\mbf{x}(t)$,
which maps a real interval $[t_0, t_1]$ into $\R ^ n$.
This is a vector-valued function
\[
\mbf{x}(t) = x_1(t)\mbf{e}_1 + x_2(t)\mbf{e}_2 + \dotsc + x_n(t)\mbf{e}_n
\]
where $\{\mbf{e}_1, \mbf{e}_2, \dotsc, \mbf{e}_n\}$ denote the Cartesian unit vectors.

If $\mbf{x}(t_1) = \mbf{x}(t_2)$ then the curve is called closed.

A curve that does not cross itself is called simple.

A parameterisation $x(t)$ is differentiable if it has a well-defined tangent vector,
\[
\frac{d\mbf{x}}{dt} = \lim_{dt \to 0}\frac{\mbf{x}(t + dt) - \mbf{x}(t)}{dt}.
\]

\subsection{Line integrals of scalar fields}
A scalar field on $\R ^ n$ is a map $f : \R ^ n \to \R$.

The line integral of a scalar field $f(\mbf{x})$ along a curve $C$ with parameterisation $\mbf{x}(t)$,
for $t \in [t_0, t_1]$,
is defined as
\[
\int_{C}f\,d\ell = \int_{t_0}^{t_1}f(\mbf{x}(t))\left|\frac{d\mbf{x}}{dt}\right|\,dt.
\]

\begin{proposition}
    The line integral of a scalar field is independent of the choice of parameterisation.

    \begin{proof}
        Suppose we have two parameterisations of the same curve:
        $\mbf{x}(t)$ for $t \in [t_0, t_1]$,
        and $\mbf{x}(\tau)$ for $\tau \in [\tau_0, \tau_1]$.
        By the chain rule,
        the tangent vectors are related by
        \[
        \frac{d\mbf{x}}{dt} = \frac{d\tau}{dt}\frac{d\mbf{x}}{d\tau} \implies \left|\frac{d\mbf{x}}{dt}\right| = \left|\frac{d\mbf{x}}{dt}\right| = \left|\frac{d\tau}{dt}\right|\left|\frac{d\mbf{x}}{d\tau}\right|.
        \]
        There are now to cases:

        Case $1$:
        $\mbf{x}(\tau_0) = \mbf{x}(t_0)$,
        meaning both parameterisation trace the curve in the same direction.
        Then
        \[
        \frac{d\tau}{dt} > 0 \implies \int_{t_0}^{t_1}f(\mbf{x}(t))\left|\frac{d\mbf{x}}{dt}\right|\,dt = \int_{t_0}^{t_1}f(\mbf{x}(t))\left|\frac{d\mbf{x}}{d\tau}\right|\frac{d\tau}{dt}\,dt = \int_{\tau_0}^{\tau_1}f(\mbf{x}(\tau))\left|\frac{d\mbf{x}}{d\tau}\right|\,d\tau.
        \]
        Case $2$:
        $\mbf{x}(\tau_0) = \mbf{x}(t_1)$,
        meaning they trace the curve in opposite directions.
        Now
        \[
        \frac{d\tau}{dt} < 0 \implies \int_{t_0}^{t_1}f(\mbf{x}(t))\left|\frac{d\mbf{x}}{dt}\right|\,dt = \int_{t_0}^{t_1}f(\mbf{x}(t))\left|\frac{d\mbf{x}}{d\tau}\right|\left(-\frac{d\tau}{dt}\right)\,dt = -\int_{\tau_1}^{\tau_0}f(\mbf{x}(\tau))\left|\frac{d\mbf{x}}{d\tau}\right|\,d\tau.
        \]
        Both cases recover the original formula.
    \end{proof}
\end{proposition}

For the special case $f(\mbf{x}) \equiv 1$,
the line integral is called the arclength of the curve,
\[
L = \int_{C}\,d\ell = \int_{t_0}^{t_1}\left|\frac{d\mbf{x}}{dt}\right|\,dt.
\]

\subsection{Line integrals of vector fields}
A vector field is a function $\mbf{f} : \R ^ n \to \R ^ n$ that assigns a vector to each point in $\R ^ n$.
For $n = 2$,
we can visualise $\mbf{f}$ as an arrow at each point in the plane.

The unit target vector of a parameterisation is
\[
\hat{\mbf{t}} = \frac{\frac{d\mbf{x}}{dt}}{\left|\frac{d\mbf{x}}{dt}\right|}
\]

An oriented curve is a curve together with a specified
(consistent)
choice of $\hvec{t}$.
We indicate this by adding an arrow to the curve.

The line integral of a vector field $\mbf{f}(\mbf{x})$ along an oriented curve $C$ with parameterisation $\vec{x}(t)$,
for $t \in [t_0, t_1]$,
is defined as the scalar line integral of the tangential component:
\[
\int_{C}\vec{f}\cdot\hvec{t}\,d\ell = \int_{t_0}^{t_1}\vec{f}(\vec{x}(t))\cdot\hvec{t}\left|\frac{d\vec{x}}{dt}\right|\,dt = \int_{t_0}^{t_1}\vec{f}(\vec{x}(t))\cdot\frac{d\vec{x}}{dt}\,dt.
\]

If the curve $C$ is closed,
this is often emphasised by the notation
\[
\int_{C}\vec{f}\cdot\hvec{t}\,d\ell = \oint_{C}\vec{f}\cdot \hvec{t}\,d\ell,
\]
and the line integral is called the circulation of $\vec{f}$ around $C$.

\subsection{Curvilinear coordinates}
In a curvilinear coordinate system $(u, v, w)$,
the coordinate lines -
where two of the coordinates are constant and one varies -
may be curved.
These coordinates lines may be parameterised by the corresponding coordinate,
so are given by $\vec{x}(u, v, w)$ with two of the coordinates fixed.

The coordinate system is orthogonal if the three tangent vectors are mutually orthogonal at every point.
We can then define an orthonormal basis $\{\mbf{e}_u, \mbf{e}_v, \mbf{e}_w\}$ by normalising:
\[
\vec{e}_u = \frac{1}{h_u}\pd[\vec{x}]{u}\quad\vec{e}_v = \frac{1}{h_v}\pd[\vec{x}]{v},\quad\vec{e}_w = \frac{1}{h_w}\pd[\vec{x}]{w},
\]
where
\[
h_u = \left|\pd[\vec{x}]{u}\right|,\quad h_v = \left|\pd[\vec{x}]{v}\right|,\quad h_w = \left|\pd[\vec{x}]{w}\right|
\]
are called the scale factors.

In $\R ^ 3$,
there are two non-Cartesian systems in common use:

\textbf{Cylindrical coordinates} are useful for problems with axial symmetry:
\[
\vec{x}(r, \theta, z) = r\cos{\theta}\vec{e}_1 + r\sin{\theta}\vec{e}_2 + z\vec{e}_3,\quad r > 0,\quad \theta \in [0, 2\pi],\quad z \in \R,
\]
leading to scale factors $h_r = 1$,
$h_{\theta} = r$,
$h_z = 1$.

\textbf{Spherical coordinates} are useful for problems with spherical symmetry:
\[
\vec{x}(r, \theta, \phi) = r\sin{\theta}\cos{\phi}\vec{e}_1 + r\sin{\theta}\sin{\phi}\vec{e}_2 + r\cos{\theta}\vec{e}_3,\quad r > 0,\quad \theta \in [0, \pi],\quad \phi \in [0, 2\pi],
\]
leading to scale factors $h_r = 1$,
$h_{\theta} = r$,
$h_{\phi} = r\sin{\theta}$.



\end{document}